%!TEX program = xelatex
% Compilar: xelatex executive_summary.tex  (dos pasadas para tikz)
\documentclass[11pt, a4paper]{article}

% ── Fuentes (XeLaTeX) ─────────────────────────────────────────────────────────
\usepackage{fontspec}
\setmainfont{Segoe UI}[
  BoldFont       = Segoe UI Bold,
  ItalicFont     = Segoe UI Italic,
  BoldItalicFont = Segoe UI Bold Italic
]
\setsansfont{Segoe UI}[
  BoldFont       = Segoe UI Bold,
  ItalicFont     = Segoe UI Italic,
  BoldItalicFont = Segoe UI Bold Italic
]
\setmonofont{Consolas}[
  BoldFont = Consolas Bold
]

% ── Geometría ─────────────────────────────────────────────────────────────────
\usepackage[
  a4paper,
  top         = 2.8cm,
  bottom      = 2.4cm,
  left        = 2.2cm,
  right       = 2.2cm,
  headheight  = 30pt,
  headsep     = 14pt
]{geometry}

% ── Colores ───────────────────────────────────────────────────────────────────
\usepackage[table]{xcolor}
\definecolor{cNavy}   {HTML}{0D1B3E}
\definecolor{cBlue}   {HTML}{2A4F9B}
\definecolor{cAccent} {HTML}{4A7FD4}
\definecolor{cBody}   {HTML}{1A1A2E}
\definecolor{cMuted}  {HTML}{6B7FA8}
\definecolor{cCodeFG} {HTML}{1B3A70}
\definecolor{cCodeBG} {HTML}{EEF2FA}
\definecolor{cTHbg}   {HTML}{0D1B3E}
\definecolor{cTRalt}  {HTML}{F2F5FC}
\definecolor{cRule}   {HTML}{C0CADF}
\definecolor{cQuoteBG}{HTML}{F0F5FF}

% ── Matemáticas ───────────────────────────────────────────────────────────────
\usepackage{amsmath}
\usepackage{amssymb}

% ── Tipografía ────────────────────────────────────────────────────────────────
\usepackage{microtype}
\usepackage{parskip}
\setlength{\parindent}{0pt}
\setlength{\parskip}{5pt plus 1pt}

% ── Tablas ────────────────────────────────────────────────────────────────────
\usepackage{booktabs}
\usepackage{tabularx}
\usepackage{colortbl}
\usepackage{array}
\usepackage{adjustbox}
\renewcommand{\arraystretch}{1.4}
\newcolumntype{L}[1]{>{\raggedright\arraybackslash}p{#1}}
\newcolumntype{C}[1]{>{\centering\arraybackslash}p{#1}}

% ── Cajas de código y citas ───────────────────────────────────────────────────
\usepackage{tcolorbox}
\tcbuselibrary{skins, breakable}

\tcbset{
  codebox/.style = {
    enhanced,
    colback    = cCodeBG,
    colframe   = cAccent,
    arc        = 3pt,
    leftrule   = 3pt,
    rightrule  = 0.5pt,
    toprule    = 0.5pt,
    bottomrule = 0.5pt,
    left       = 10pt,
    right      = 10pt,
    top        = 7pt,
    bottom     = 7pt,
    fontupper  = \small\ttfamily\color{cCodeFG},
    breakable
  },
  quotebox/.style = {
    enhanced,
    colback    = cQuoteBG,
    colframe   = cAccent,
    arc        = 2pt,
    leftrule   = 3pt,
    rightrule  = 0pt,
    toprule    = 0pt,
    bottomrule = 0pt,
    left       = 12pt,
    right      = 10pt,
    top        = 7pt,
    bottom     = 7pt,
    fontupper  = \itshape\color{cCodeFG},
    breakable
  }
}

% ── Encabezados de sección ────────────────────────────────────────────────────
\usepackage{titlesec}

\titleformat{\section}
  {\color{cNavy}\Large\bfseries}
  {}{0em}{}
  [\vspace{2pt}{\color{cAccent}\hrule height 1.6pt}\vspace{4pt}]

\titleformat{\subsection}
  {\color{cBlue}\large\bfseries}
  {}{0em}{}
  [\vspace{1pt}{\color{cRule}\hrule height 0.5pt}\vspace{2pt}]

\titleformat{\subsubsection}
  {\color{cBlue}\normalsize\bfseries}
  {}{0em}{}

\titlespacing*{\section}      {0pt}{20pt}{8pt}
\titlespacing*{\subsection}   {0pt}{14pt}{4pt}
\titlespacing*{\subsubsection}{0pt}{10pt}{3pt}

% ── Cabecera y pie de página ──────────────────────────────────────────────────
\usepackage{fancyhdr}
\usepackage{tikz}
\usetikzlibrary{calc}

\pagestyle{fancy}
\fancyhf{}
\renewcommand{\headrulewidth}{0pt}

\fancyhead[C]{%
  \begin{tikzpicture}[remember picture, overlay]
    \fill[cNavy]
      (current page.north west)
      rectangle ($(current page.north east) + (0, -1.15cm)$);
    \fill[cAccent]
      (current page.north west)
      rectangle ($(current page.north west) + (0.28\paperwidth, -1.15cm)$);
    \node[white, anchor=west, font=\small\itshape] at
      ($(current page.north west) + (0.28\paperwidth + 8pt, -0.58cm)$)
      {Go Ising \enspace\textbullet\enspace Experimento 06
       \enspace\textbullet\enspace Familia de Hamiltonianos cúbicos};
  \end{tikzpicture}%
}

\fancyfoot[L]{\small\color{cMuted} Leonardo Jiménez Martínez \enspace·\enspace UNAM \enspace·\enspace 2026}
\fancyfoot[R]{\small\color{cMuted} Página \thepage}
\renewcommand{\footrulewidth}{0.4pt}
\def\footrule{{\color{cRule}\vspace{-3pt}\hrule width\headwidth height 0.4pt\vspace{2pt}}}

% ── Listas ────────────────────────────────────────────────────────────────────
\usepackage{enumitem}
\setlist[itemize]{
  leftmargin = 1.5em,
  itemsep    = 2pt,
  topsep     = 4pt,
  label      = {\color{cAccent}\small$\bullet$}
}
\setlist[enumerate]{leftmargin=1.8em, itemsep=2pt, topsep=4pt}

% ── PDF metadata ──────────────────────────────────────────────────────────────
\usepackage[
  colorlinks = true,
  linkcolor  = cAccent,
  urlcolor   = cAccent,
  pdfauthor  = {Leonardo Jiménez Martínez},
  pdftitle   = {La familia de polinomios Go y sus variedades asociadas},
  pdfsubject = {Go Ising · Experimento 06 · Hamiltonianos cúbicos}
]{hyperref}

% ── Macros de tabla ───────────────────────────────────────────────────────────
\newcommand{\TH}[1]{%
  \cellcolor{cTHbg}\color{white}\textbf{#1}%
}

% =============================================================================
\begin{document}
\thispagestyle{fancy}

% ── Portada ───────────────────────────────────────────────────────────────────
\vspace*{4pt}
{\color{cNavy}\Huge\bfseries La familia de polinomios Go}\\[3pt]
{\color{cNavy}\Large\bfseries y sus variedades asociadas}

\vspace{6pt}
{\color{cMuted}\large\itshape
  Informe ejecutivo --- Fibración de Milnor en estrategia de Go}

\vspace{4pt}
{\color{cMuted}\small\itshape Generado automáticamente · 2026-07-15}\\[1pt]
{\color{cMuted}\small\itshape
  Catálogo: 305 Hamiltonianos (5 referencias + 300 candidatos) · 296 pasan el filtro}

\vspace{8pt}
{\color{cAccent}\hrule height 2pt}
\vspace{10pt}

% =============================================================================
\section{La familia Go de Hamiltonianos cúbicos}

\subsection{El modelo de espín en Go}

En el modelo de espín de Go, cada intersección del tablero tiene un valor
$s \in \{-1,\, 0,\, +1\}$:

\medskip
\begin{center}
\rowcolors{2}{white}{cTRalt}
\begin{tabular}{cc}
  \TH{Espín} & \TH{Significado} \\[1pt]
  $s = -1$ & Piedra negra \\
  $s = \phantom{-}0$  & Intersección vacía \\
  $s = +1$ & Piedra blanca \\
\end{tabular}
\end{center}
\medskip

Un \textbf{Hamiltoniano de interacción} $H(s_i, s_j)$ asigna una energía a cada par
de intersecciones vecinas $(i, j)$. La energía total del tablero en una configuración
dada es la suma sobre todos los pares adyacentes:

\begin{tcolorbox}[codebox]
$\displaystyle E \;=\; \sum_{\langle i,j \rangle} H(s_i,\; s_j)$
\end{tcolorbox}

Solo existen 6 pares orientados relevantes (excluimos la dirección inversa como
independiente):

\begin{tcolorbox}[codebox]
$(-1,-1)$ \quad $(-1,0)$ \quad $(-1,1)$ \quad $(1,-1)$ \quad $(1,0)$ \quad $(1,1)$
\end{tcolorbox}

Estos 6 valores determinan completamente el comportamiento estratégico del modelo:
cómo penaliza o favorece cada tipo de vecindad negro--negro, negro--vacío,
negro--blanco, etc.

% ─────────────────────────────────────────────────────────────────────────────
\subsection{¿Por qué polinomios? La base canónica en $\{-1,\,0,\,+1\}$}

Sobre el conjunto discreto $\{-1, 0, +1\}$, todo polinomio en $x$ satisface
la identidad:

\begin{tcolorbox}[codebox]
$x^3 = x \qquad \forall\; x \in \{-1,\, 0,\, +1\}$
\end{tcolorbox}

\noindent
(verificar: $(-1)^3 = -1$, $\; 0^3 = 0$, $\; 1^3 = 1$). Esto implica que los
monomios de grado $\geq 3$ en una sola variable son linealmente dependientes de
los de menor grado. Los monomios \textbf{independientes} en
$(x,y) \in \{-1,0,+1\}^2$ son exactamente:

\begin{tcolorbox}[codebox]
$1, \quad x, \quad y, \quad x^2, \quad xy, \quad y^2, \quad x^2 y, \quad xy^2$
\end{tcolorbox}

Descartando el término constante (desplaza la energía globalmente sin efecto
estratégico relativo), obtenemos \textbf{7 monomios base}. Cualquier función
$H\colon \{-1,0,+1\}^2 \to \mathbb{R}$ puede expresarse como combinación
lineal de estos 7 términos. Los polinomios cúbicos son, en este sentido,
\textbf{la familia más general posible} de Hamiltonianos de interacción por par
para Go.

% ─────────────────────────────────────────────────────────────────────────────
\subsection{La plantilla \texttt{cubic\_mixed}: el espacio completo de 7 parámetros}

\begin{tcolorbox}[codebox]
$H(x,y) \;=\; a_1 x \;+\; a_2 y \;+\; b_{11} x^2 \;+\; b_{12}\, xy \;+\;
b_{22} y^2 \;+\; c_{112}\, x^2 y \;+\; c_{122}\, x y^2$
\end{tcolorbox}

Cada coeficiente controla una dimensión independiente de la interacción:

\medskip
\begin{center}
\rowcolors{2}{white}{cTRalt}
\begin{tabularx}{0.96\textwidth}{C{2cm}C{2cm}X}
  \TH{Parámetro} & \TH{Monomio} & \TH{Rol estratégico} \\[1pt]
  $a_1$     & $x$       & Tendencia lineal de la piedra $i$
                           (favorece negro o blanco puro) \\
  $a_2$     & $y$       & Tendencia lineal de la piedra $j$ \\
  $b_{11}$  & $x^2$     & Energía cuadrática de $i$ --- igual para
                           $-1$ y $+1$ (simétrica en color) \\
  $b_{12}$  & $xy$      & Interacción directa $i$--$j$ ---
                           el núcleo Ising clásico \\
  $b_{22}$  & $y^2$     & Energía cuadrática de $j$ \\
  $c_{112}$ & $x^2 y$   & Interacción cruzada: el color de $j$
                           modulado por si $i$ es vacío o no \\
  $c_{122}$ & $x y^2$   & Interacción cruzada: el color de $i$
                           modulado por si $j$ es vacío o no \\
\end{tabularx}
\end{center}
\medskip

Los términos $c_{112}$ y $c_{122}$ son los que generan \textbf{asimetría
direccional}: distinguen si la piedra relevante está en la posición $i$ o en $j$
del par orientado. El modelo $H_{M1}$ de Mercado \& Jiménez usa exactamente este
mecanismo con $b_{12}=0$, $a_1=1$, $a_2=2$, $c_{112}=-1$, $c_{122}=-1$.

% ─────────────────────────────────────────────────────────────────────────────
\subsection{Los Hamiltonianos de referencia}

Tres miembros especiales fijan el marco de comparación:

\medskip
\begin{center}
\rowcolors{2}{white}{cTRalt}
\begin{tabularx}{0.96\textwidth}{L{2cm}L{2.4cm}L{4.4cm}X}
  \TH{ID} & \TH{Plantilla} & \TH{Expresión} & \TH{Por qué es referencia} \\[1pt]
  $H_{M1}$
    & \texttt{sparse\_cubic}
    & $x + 2y - xy^2 - x^2y$
    & Modelo asimétrico con vacío activo; base del paper de josekis
      (Jiménez \& Sesma, 2025) \\
  Alvarado
    & \texttt{quadratic}
    & $xy$
    & Modelo simétrico canónico con vacío invisible; Atomic-Go
      (Rojas-Domínguez et al., 2019) \\
  Armónico
    & \texttt{sym\_cubic}
    & $x^3+y^3 - 3xy - 3(x^2y+xy^2)$
    & Función armónica ($\Delta f = 0$) y simétrica $f(x,y)=f(y,x)$;
      2 nodos $A_1$ \\
\end{tabularx}
\end{center}
\medskip

% ─────────────────────────────────────────────────────────────────────────────
\subsection{¿Por qué 300 muestras aleatorias?}

El espacio de coeficientes de \texttt{cubic\_mixed} es $\mathbb{R}^7$. Un
muestreo en cuadrícula con solo 3 valores por eje daría $3^7 = 2187$
combinaciones, con alta redundancia. El muestreo aleatorio uniforme cubre el
espacio de forma más eficiente:

\begin{itemize}
  \item $a_1,\; a_2,\; b_{11},\; b_{12},\; b_{22} \;\in\; [-3,\; 3]$
        \hfill\textit{(términos dominantes)}
  \item $c_{112},\; c_{122} \;\in\; [-1,\; 1]$
        \hfill\textit{(términos cúbicos --- perturbación del núcleo cuadrático)}
\end{itemize}

Con semilla fija (\texttt{seed = 42}) el muestreo es completamente reproducible.
Los rangos reflejan que los términos cúbicos actúan como correcciones de tercer
orden, no como el término dominante de la energía.

% =============================================================================
\section{Resumen de la búsqueda}

Plantillas exploradas:

\begin{itemize}
  \item \texttt{cubic\_mixed}: 300 muestras aleatorias
  \item \texttt{h\_m1}: 2 muestras (referencia $H_{M1}$)
  \item \texttt{quadratic}: 2 muestras (referencia Alvarado)
  \item \texttt{sym\_cubic}: 1 muestra (referencia armónica)
\end{itemize}

\medskip
\begin{center}
\rowcolors{2}{white}{cTRalt}
\begin{tabular}{lr}
  \TH{Métrica} & \TH{Valor} \\[1pt]
  Total analizados  & 305 \\
  Pasan el filtro   & 296 \\
  Tasa de éxito     & $97.0\,\%$ \\
  Top candidatos    & 5 \\
\end{tabular}
\end{center}

% =============================================================================
\section{Top 5 candidatos}

\medskip
\begin{center}
\footnotesize
\begin{adjustbox}{max width=\textwidth}
\rowcolors{2}{white}{cTRalt}
\begin{tabular}{llp{6.5cm}C{0.7cm}C{1.1cm}C{0.9cm}C{0.9cm}C{1.1cm}C{0.9cm}r}
  \TH{ID} & \TH{Template} & \TH{Expresión} &
  \TH{$A_1$} & \TH{$\Delta c$} & \TH{$\Delta E$} &
  \TH{Rob.} & \TH{Mejora} & \TH{$p$-val} & \TH{Score} \\[1pt]
  H\_0094 & cubic\_mixed & \texttt{-0.781435\,x\textasciicircum2*y - 2.778\,...}
    & 1 & 9.914  & 31.3 & 1.000 & +0.0000 & --- & \textbf{0.4593} \\
  H\_0022 & cubic\_mixed & \texttt{+0.850240\,x\textasciicircum2*y - 2.380\,...}
    & 1 & 16.779 & 23.5 & 1.000 & +0.0000 & --- & \textbf{0.4575} \\
  H\_0113 & cubic\_mixed & \texttt{-0.778741\,x\textasciicircum2*y - 2.944\,...}
    & 1 & 22.326 & 35.4 & 1.000 & +0.0000 & --- & \textbf{0.4559} \\
  H\_0045 & cubic\_mixed & \texttt{+0.849684\,x\textasciicircum2*y - 2.390\,...}
    & 2 & 13.381 & 43.4 & 1.000 & +0.0000 & --- & \textbf{0.4544} \\
  H\_0275 & cubic\_mixed & \texttt{+0.697584\,x\textasciicircum2*y - 2.670\,...}
    & 1 & 20.716 & 33.5 & 1.000 & +0.0000 & --- & \textbf{0.4540} \\
\end{tabular}
\end{adjustbox}
\end{center}

\smallskip
\begin{tcolorbox}[quotebox]
  \textbf{$A_1$}: nodos $A_1$ detectados \;\textbullet\;
  \textbf{$\Delta c$}: separación mínima entre valores críticos $c^*$ \;\textbullet\;
  \textbf{$\Delta E$}: rango energético en $[-2,2]^2$ \;\textbullet\;
  \textbf{Rob.}: fracción de perturbaciones $\pm 5\%$ estables
\end{tcolorbox}

% =============================================================================
\section{Criterios de selección}

Un candidato pasa si cumple \textbf{al menos 2 de 3} de los siguientes criterios:

\medskip
\begin{center}
\rowcolors{2}{white}{cTRalt}
\begin{tabularx}{0.96\textwidth}{L{3.2cm}C{2cm}X}
  \TH{Criterio} & \TH{Umbral} & \TH{Qué mide} \\[1pt]
  $H_1^{\max}$ ($\tau_1$)      & $> 0.20$ &
    Loops topológicos en subniveles de $H$ ---
    complejidad táctica de la función de evaluación \\
  $\Delta E$                   & $> 0.50$ &
    Rango energético (profundidad del pozo) ---
    contraste entre posiciones extremas \\
  $\delta_{\mathrm{crit}}$     & $> 0.30$ &
    Separación mínima entre valores críticos $c^*$ ---
    resolución entre regímenes estratégicos \\
\end{tabularx}
\end{center}

\medskip
\textbf{Score total}
$\;=\; 0.45 \times \mathrm{TDA} \;+\; 0.30 \times \mathrm{Robustez}
        \;+\; 0.25 \times \mathrm{Mejora\; estratégica}$

% =============================================================================
\section{Invariantes matemáticos}

\medskip
\begin{center}
\rowcolors{2}{white}{cTRalt}
\begin{tabularx}{0.96\textwidth}{L{3cm}L{5.5cm}X}
  \TH{Invariante} & \TH{Definición} & \TH{Rol en Go} \\[1pt]
  Nodo $A_1$
    & $\det(\mathrm{Hess}\,H) < 0$ en punto crítico
    & Fibra singular $\to$ tensión crítica en la partida \\
  Número de Milnor $\mu$
    & $\#$ nodos $A_1$;\; $\mu \leq (d-1)^2$
    & Complejidad topológica de la fibración \\
  Genus $g$ de fibra
    & $g=(d-1)(d-2)/2$;\; cúbico $\Rightarrow g=1$\,(toro)
    & Tipo topológico de cada estado de la partida \\
  $c^*$ (valor crítico)
    & $H(\text{punto crítico})$
    & Umbral de cambio topológico \\
  Separación $\Delta c$
    & $\min_i |c^*_i - c^*_j|$
    & Resolución entre regímenes estratégicos \\
\end{tabularx}
\end{center}

\vfill
\noindent
{\color{cRule}\hrule height 0.5pt}
\vspace{4pt}
{\color{cMuted}\small\itshape
  Pipeline: \texttt{experiments/06\_hamiltonian\_families/} \quad\textbullet\quad
  Catálogo: \texttt{output/catalog.json}
}

\end{document}
